\paragraph{Description}
% The sampling region where the sample is located: one of 200 regions of County and PUMA cells with similar heating and cooling climate, each within one state and one UTC offset.

The sampling region where the sample is located. Sampling regions group County and PUMA cells of similar heating and cooling climate, each within one state and one UTC offset. They are the geographic segment variable of the stratified sampler, and the allocation only matches a building to housing units in its own sampling region. Section \ref{sec:sampling_region_characteristic} describes how the regions are built.

\paragraph{Distribution Data Sources}
% \begin{itemize}
\item Clustered on the AMY2018 county weather file assigned to each County and PUMA cell, weighted by ACS 2015-2019 5-year housing unit counts.
\end{itemize}

\begin{itemize}
    \item Regions are clustered on the AMY2018 county weather file assigned to each County and PUMA cell, weighted by ACS 2015--2019 5-year housing unit counts.
    \item Spatial definitions are from the U.S.~Census Bureau 2010 county, PUMA and tract boundaries.
\end{itemize}

\paragraph{Direct Conditional Dependencies}
\begin{itemize}
\item CEC Climate Zone
\item County
\end{itemize}


\paragraph{Options}
200 regions, each named by a numeric identifier: state FIPS code $\times$ 10,000 + $|$UTC offset$|$ $\times$ 100 + $k$, where $k$ ranks the region by its heating and cooling energy index within its state and UTC offset, 1 being the mildest. For example, region 10601 is the mildest region of Alabama (FIPS 01) on UTC$-$6.

\paragraph{Distribution Assumptions}
% \begin{itemize}
\item Every County and PUMA cell belongs to exactly one region.
\end{itemize}

Every County and PUMA cell belongs to exactly one region, so the distribution is deterministic given County and PUMA.
